\chapter{Cấu hình và protocol tái lập}\label{app:config}
\begingroup\small\setlength{\tabcolsep}{4pt}\renewcommand{\arraystretch}{1.16}
\begin{longtable}{p{.19\linewidth}p{.40\linewidth}p{.31\linewidth}}
\caption{Cấu hình thực nghiệm và app hiện tại; phạm vi được phân biệt theo thành phần.}\label{tab:configuration}\\
\toprule Thành phần & Tham số & Giá trị \\
\midrule\endfirsthead
\multicolumn{3}{l}{\textit{Bảng tiếp theo}}\\
\toprule Thành phần & Tham số & Giá trị \\
\midrule\endhead
\bottomrule\endfoot
Stage 1 & Language confidence & 0,90 \\
Stage 1 & Title / Abstract near-duplicate thresholds & 0,88 / 0,80 \\
LDA & K khảo sát / passes / iterations & 10,15,20,25,30 / 10 / 100 \\
LDA & alpha / eta / seed cơ sở & asymmetric / symmetric / 42 \\
BGE & Dimension / train max tokens & 768 / 128 \\
Fine-tuning & Epoch / batch / learning rate & 1 / 8 / $2\times10^{-5}$ \\
Fine-tuning & Block trainable / temperature & 1 block cuối / 0,05 \\
Chunking & Max pair / overlap / pooling & 256 / 64 / normalized mean CLS \\
UMAP & Neighbors / components / min-dist & 15 / 5 / 0 \\
HDBSCAN & Min cluster / min samples / selection & 30 / 10 / EOM \\
Search v2 & Candidate / rerank depth / RRF k & 100 / 50 / 60 \\
Search v2 & Title weight / BM25 k1 / b & 2 / 1,5 / 0,75 \\
Reranker & Passage / overlap / max passages & 320 / 64 / 2 \\
Reranker & Max pair / Title budget / batch & 512 / 80 / 4 \\
Gemini & Model A/B/C/generator & gemini-3.5-flash-lite \\
Evaluation & Bootstrap repeats / seed & 2.000 / 42 \\
\end{longtable}\endgroup
\source{configs/config.yaml; configs/stage2.yaml; configs/stage3.yaml; configs/search_v2.json; outputs/retrieval_ai_evaluation_v1/protocol.json}

Cấu hình trong Bảng~\ref{tab:configuration} được đọc từ files hiện tại và đối chiếu các contracts của artifact. Việc mã nguồn hiện tại chứa alias Stage4 không có nghĩa implementation cũ đã bị mất: ZIP archive và manifest giữ nguồn trước cleanup. Công thức SE1/SE2/SE3 trong báo cáo dựa vào bản archived, trong khi mô tả app dựa vào Search v2 hiện hành.

\section{Quy tắc diễn giải trạng thái}
\artifact{AI_CONSENSUS_UNVERIFIED} chỉ nhãn AI đã qua quy tắc chấp nhận; \artifact{PENDING_JUDGES} thiếu lượt chấm; \artifact{PENDING_ADJUDICATION} đang chờ C; UNRESOLVED là chưa đủ căn cứ. Không cộng các nhóm chưa đủ điều kiện vào tử số hay mẫu số relevance. Chỉ Q11/Q19 có full pool tại snapshot, nên hai query được dùng chung cho mọi method; development/remaining lần lượt mới có 1/5 và 1/18.

Sau khi có nhãn bổ sung, phải sinh lại các bảng theo cùng protocol và đọc lại phần nhận xét, không chỉ thay số rồi giữ nguyên các kết luận partial. Script dựng bảng từ chối snapshot đã đổi khỏi điều kiện biên soạn hiện tại để tránh âm thầm làm nội dung tĩnh sai. API key và nội dung tệp môi trường không được script báo cáo đọc hoặc sao chép.

\section{Kiểm chứng và đơn vị số liệu}
Mọi mean relevance là macro-average theo query đủ pool; synthetic là exact source-ID recovery theo query đã chấp nhận. Kết quả LDA/BERTopic dùng keywords theo topic và các tập noise khác nhau. Khoảng tin cậy có điều kiện trên các nhãn/checkpoints cố định, không biến hai query thành một mẫu lớn hơn. Giá trị phần trăm noise sử dụng denominator 13.702; điểm tăng noise được mô tả theo điểm phần trăm, khác tỷ lệ tăng tương đối.

\chapter{Truy vấn và kết quả mở rộng}\label{app:queries}
\begingroup\small\setlength{\tabcolsep}{4pt}\renewcommand{\arraystretch}{1.16}
\begin{longtable}{lp{.67\linewidth}l}
\caption{Truy vấn miền đã lưu; nội dung tiếng Anh giữ nguyên.}\label{tab:queries}\\
\toprule ID & Query & Vai trò \\
\midrule\endfirsthead
\multicolumn{3}{l}{\textit{Bảng tiếp theo}}\\
\toprule ID & Query & Vai trò \\
\midrule\endhead
\bottomrule\endfoot
Q01 & Hard carbon anodes derived from biomass for sodium-ion batteries & Dev \\
Q02 & Initial Coulombic efficiency and sodium storage mechanisms in hard carbon & Authored \\
Q03 & P2 and O3 layered oxide cathodes with reversible anionic redox & Authored \\
Q04 & Prussian blue analogues with low vacancies and low water content & Authored \\
Q05 & Na3V2(PO4)3 NASICON cathodes with carbon coating & Authored \\
Q06 & Low-cost iron phosphate and pyrophosphate cathode materials & Authored \\
Q07 & Fluoroethylene carbonate additives and stable SEI on hard carbon & Authored \\
Q08 & NaFSI ether electrolytes for low-temperature sodium-ion batteries & Dev \\
Q09 & Na3PS4 and Na3SbS4 sulfide solid electrolytes & Authored \\
Q10 & NASICON Na3Zr2Si2PO12 ceramic ionic conductivity & Authored \\
Q11 & Ti3C2Tx MXene interlayer engineering for sodium storage & Authored \\
Q12 & MoS2 nanosheets and expanded interlayer spacing for sodium storage & Authored \\
Q13 & Sn4P3 alloying anodes and volume expansion control & Authored \\
Q14 & Organic quinone and covalent framework sodium-ion electrodes & Authored \\
Q15 & Slurry binders and current collectors for sodium-ion electrodes & Authored \\
Q16 & Roll pressing and electrode coating equipment for sodium-ion batteries & Dev \\
Q17 & Thermal management and safety of sodium-ion battery packs & Authored \\
Q18 & Recycling spent sodium-ion batteries and recovery of cathode metals & Authored \\
Q19 & Aqueous sodium-ion batteries with water-in-salt electrolytes & Dev \\
Q20 & Hard carbon full cells with layered oxide cathodes and electrolyte additives & Authored \\
Q21 & Operando characterization of sodium intercalation and phase transitions & Authored \\
Q22 & Cellulose separators with improved electrolyte wettability & Authored \\
Q23 & Na2FeP2O7 cathode electrochemical cycling stability & Dev \\
\end{longtable}\endgroup
\source{outputs/tables/stage4/queries.csv; outputs/retrieval_ai_evaluation_v1/protocol.json}

Q24 là truy vấn OOV control, được lưu riêng và loại khỏi mean miền. Nội dung query xuất từ artifact; đánh dấu development thể hiện đã dùng để phát triển, không khẳng định 18 query còn lại là test ngoài độc lập.
\begingroup\small\setlength{\tabcolsep}{4pt}\renewcommand{\arraystretch}{1.16}
\begin{longtable}{lrrl}
\caption{Coverage nhãn từng query; chỉ query đủ pool vào mean chung.}\label{tab:coverage}\\
\toprule Query & Pool & Resolved & Đưa vào metrics \\
\midrule\endfirsthead
\multicolumn{4}{l}{\textit{Bảng tiếp theo}}\\
\toprule Query & Pool & Resolved & Đưa vào metrics \\
\midrule\endhead
\bottomrule\endfoot
Q01 & 43 & 0 & Chưa \\
Q02 & 40 & 0 & Chưa \\
Q03 & 38 & 0 & Chưa \\
Q04 & 49 & 0 & Chưa \\
Q05 & 44 & 0 & Chưa \\
Q06 & 44 & 0 & Chưa \\
Q07 & 41 & 0 & Chưa \\
Q08 & 49 & 0 & Chưa \\
Q09 & 26 & 0 & Chưa \\
Q10 & 37 & 0 & Chưa \\
Q11 & 37 & 37 & Có \\
Q12 & 47 & 0 & Chưa \\
Q13 & 41 & 0 & Chưa \\
Q14 & 34 & 0 & Chưa \\
Q15 & 47 & 0 & Chưa \\
Q16 & 44 & 0 & Chưa \\
Q17 & 37 & 0 & Chưa \\
Q18 & 38 & 36 & Chưa \\
Q19 & 43 & 43 & Có \\
Q20 & 48 & 0 & Chưa \\
Q21 & 53 & 0 & Chưa \\
Q22 & 29 & 0 & Chưa \\
Q23 & 33 & 0 & Chưa \\
\end{longtable}\endgroup
\source{outputs/retrieval_ai_evaluation_v1/query_coverage.csv}

\begingroup\small\setlength{\tabcolsep}{4pt}\renewcommand{\arraystretch}{1.16}
\begin{longtable}{p{.34\linewidth}lrrr}
\caption{Source recovery theo nguồn; 97 Patent và 100 Publication.}\label{tab:synthetic-source}\\
\toprule Method & Nhóm & $n$ & Hit@10 & MRR@10 \\
\midrule\endfirsthead
\multicolumn{5}{l}{\textit{Bảng tiếp theo}}\\
\toprule Method & Nhóm & $n$ & Hit@10 & MRR@10 \\
\midrule\endhead
\bottomrule\endfoot
{\footnotesize\texttt{BM25}} & Patent & 97 & 0,9588 & 0,8650 \\
{\footnotesize\texttt{BM25}} & Publication & 100 & 0,8400 & 0,7355 \\
{\footnotesize\texttt{BM25\_CHEMICAL}} & Patent & 97 & 0,9691 & 0,8172 \\
{\footnotesize\texttt{BM25\_CHEMICAL}} & Publication & 100 & 0,7100 & 0,6020 \\
{\footnotesize\texttt{DENSE\_CHUNKED}} & Patent & 97 & 0,9794 & 0,9029 \\
{\footnotesize\texttt{DENSE\_CHUNKED}} & Publication & 100 & 0,8700 & 0,7860 \\
{\footnotesize\texttt{DENSE\_FINETUNED}} & Patent & 97 & 0,9381 & 0,8578 \\
{\footnotesize\texttt{DENSE\_FINETUNED}} & Publication & 100 & 0,8500 & 0,7569 \\
{\footnotesize\texttt{DENSE\_PRETRAINED}} & Patent & 97 & 0,8763 & 0,7682 \\
{\footnotesize\texttt{DENSE\_PRETRAINED}} & Publication & 100 & 0,8700 & 0,7520 \\
{\footnotesize\texttt{HYBRID\_CHUNKED}} & Patent & 97 & 0,9897 & 0,9158 \\
{\footnotesize\texttt{HYBRID\_CHUNKED}} & Publication & 100 & 0,8400 & 0,6988 \\
{\footnotesize\texttt{HYBRID\_FINETUNED}} & Patent & 97 & 0,9691 & 0,8904 \\
{\footnotesize\texttt{HYBRID\_FINETUNED}} & Publication & 100 & 0,8100 & 0,7020 \\
{\footnotesize\texttt{HYBRID\_PRETRAINED}} & Patent & 97 & 0,9485 & 0,8390 \\
{\footnotesize\texttt{HYBRID\_PRETRAINED}} & Publication & 100 & 0,8200 & 0,7112 \\
{\footnotesize\texttt{RERANK\_CHUNKED}} & Patent & 97 & 0,9485 & 0,8131 \\
{\footnotesize\texttt{RERANK\_CHUNKED}} & Publication & 100 & 0,8500 & 0,7194 \\
{\footnotesize\texttt{RERANK\_FINETUNED}} & Patent & 97 & 0,9588 & 0,8159 \\
{\footnotesize\texttt{RERANK\_FINETUNED}} & Publication & 100 & 0,8500 & 0,7174 \\
{\footnotesize\texttt{RERANK\_PRETRAINED}} & Patent & 97 & 0,9588 & 0,8136 \\
{\footnotesize\texttt{RERANK\_PRETRAINED}} & Publication & 100 & 0,8500 & 0,7181 \\
{\footnotesize\texttt{SE1}} & Patent & 97 & 0,0619 & 0,0228 \\
{\footnotesize\texttt{SE1}} & Publication & 100 & 0,0500 & 0,0189 \\
{\footnotesize\texttt{SE2}} & Patent & 97 & 0,9381 & 0,8359 \\
{\footnotesize\texttt{SE2}} & Publication & 100 & 0,8500 & 0,7606 \\
{\footnotesize\texttt{SE3}} & Patent & 97 & 0,7835 & 0,3411 \\
{\footnotesize\texttt{SE3}} & Publication & 100 & 0,7800 & 0,3178 \\
\end{longtable}\endgroup
\source{outputs/retrieval_ai_evaluation_v1/synthetic/by_source.csv}

\begingroup\setstretch{1.0}\footnotesize\setlength{\tabcolsep}{4pt}\renewcommand{\arraystretch}{1.10}
\begin{longtable}{p{.34\linewidth}lrrr}
\caption{Source recovery theo split của source document; không phải relevance.}\label{tab:synthetic-split}\\
\toprule Method & Nhóm & $n$ & Hit@10 & MRR@10 \\
\midrule\endfirsthead
\multicolumn{5}{l}{\textit{Bảng tiếp theo}}\\
\toprule Method & Nhóm & $n$ & Hit@10 & MRR@10 \\
\midrule\endhead
\bottomrule\endfoot
{\footnotesize\texttt{BM25}} & test & 18 & 0,8889 & 0,7625 \\
{\footnotesize\texttt{BM25}} & train & 163 & 0,9018 & 0,8175 \\
{\footnotesize\texttt{BM25}} & validation & 16 & 0,8750 & 0,6547 \\
{\footnotesize\texttt{BM25\_CHEMICAL}} & test & 18 & 0,8333 & 0,7346 \\
{\footnotesize\texttt{BM25\_CHEMICAL}} & train & 163 & 0,8405 & 0,7171 \\
{\footnotesize\texttt{BM25\_CHEMICAL}} & validation & 16 & 0,8125 & 0,5854 \\
{\footnotesize\texttt{DENSE\_CHUNKED}} & test & 18 & 0,9444 & 0,8750 \\
{\footnotesize\texttt{DENSE\_CHUNKED}} & train & 163 & 0,9202 & 0,8385 \\
{\footnotesize\texttt{DENSE\_CHUNKED}} & validation & 16 & 0,9375 & 0,8594 \\
{\footnotesize\texttt{DENSE\_FINETUNED}} & test & 18 & 0,8889 & 0,8611 \\
{\footnotesize\texttt{DENSE\_FINETUNED}} & train & 163 & 0,8896 & 0,8040 \\
{\footnotesize\texttt{DENSE\_FINETUNED}} & validation & 16 & 0,9375 & 0,7714 \\
{\footnotesize\texttt{DENSE\_PRETRAINED}} & test & 18 & 0,8889 & 0,8009 \\
{\footnotesize\texttt{DENSE\_PRETRAINED}} & train & 163 & 0,8712 & 0,7503 \\
{\footnotesize\texttt{DENSE\_PRETRAINED}} & validation & 16 & 0,8750 & 0,8125 \\
{\footnotesize\texttt{HYBRID\_CHUNKED}} & test & 18 & 0,8889 & 0,7963 \\
{\footnotesize\texttt{HYBRID\_CHUNKED}} & train & 163 & 0,9141 & 0,8078 \\
{\footnotesize\texttt{HYBRID\_CHUNKED}} & validation & 16 & 0,9375 & 0,7937 \\
{\footnotesize\texttt{HYBRID\_FINETUNED}} & test & 18 & 0,8889 & 0,7824 \\
{\footnotesize\texttt{HYBRID\_FINETUNED}} & train & 163 & 0,8896 & 0,7995 \\
{\footnotesize\texttt{HYBRID\_FINETUNED}} & validation & 16 & 0,8750 & 0,7604 \\
{\footnotesize\texttt{HYBRID\_PRETRAINED}} & test & 18 & 0,8889 & 0,7370 \\
{\footnotesize\texttt{HYBRID\_PRETRAINED}} & train & 163 & 0,8896 & 0,7816 \\
{\footnotesize\texttt{HYBRID\_PRETRAINED}} & validation & 16 & 0,8125 & 0,7396 \\
{\footnotesize\texttt{RERANK\_CHUNKED}} & test & 18 & 0,8889 & 0,7593 \\
{\footnotesize\texttt{RERANK\_CHUNKED}} & train & 163 & 0,8896 & 0,7653 \\
{\footnotesize\texttt{RERANK\_CHUNKED}} & validation & 16 & 1,0000 & 0,7746 \\
{\footnotesize\texttt{RERANK\_FINETUNED}} & test & 18 & 0,8889 & 0,7500 \\
{\footnotesize\texttt{RERANK\_FINETUNED}} & train & 163 & 0,8957 & 0,7668 \\
{\footnotesize\texttt{RERANK\_FINETUNED}} & validation & 16 & 1,0000 & 0,7746 \\
{\footnotesize\texttt{RERANK\_PRETRAINED}} & test & 18 & 0,8889 & 0,7500 \\
{\footnotesize\texttt{RERANK\_PRETRAINED}} & train & 163 & 0,8957 & 0,7659 \\
{\footnotesize\texttt{RERANK\_PRETRAINED}} & validation & 16 & 1,0000 & 0,7746 \\
{\footnotesize\texttt{SE1}} & test & 18 & 0,0556 & 0,0079 \\
{\footnotesize\texttt{SE1}} & train & 163 & 0,0552 & 0,0182 \\
{\footnotesize\texttt{SE1}} & validation & 16 & 0,0625 & 0,0625 \\
{\footnotesize\texttt{SE2}} & test & 18 & 0,8889 & 0,8194 \\
{\footnotesize\texttt{SE2}} & train & 163 & 0,8896 & 0,7981 \\
{\footnotesize\texttt{SE2}} & validation & 16 & 0,9375 & 0,7688 \\
{\footnotesize\texttt{SE3}} & test & 18 & 0,8889 & 0,4224 \\
{\footnotesize\texttt{SE3}} & train & 163 & 0,7853 & 0,3243 \\
{\footnotesize\texttt{SE3}} & validation & 16 & 0,6250 & 0,2751 \\
\end{longtable}\endgroup
\source{outputs/retrieval_ai_evaluation_v1/synthetic/by_training_split.csv}

\chapter{Bản đồ bằng chứng và khả năng tái lập}\label{app:evidence}
Thư mục báo cáo chứa bản sao các hình đã sử dụng, không sửa biểu đồ thực nghiệm gốc. \artifact{evidence_manifest.json} lưu SHA-256 nguồn và phép tổng hợp bảng; \artifact{code_review_inventory.json} ghi các module/config được đọc để đối chiếu, không tuyên bố kiểm chứng toàn bộ hành vi runtime chỉ qua static inspection. \artifact{validation.json} ghi kiểm tra liên kết, citations, số liệu, snapshot và biên dịch khi có môi trường phù hợp.

Để tái dựng bảng, từ gốc dự án chạy \artifact{.venv/Scripts/python.exe report_latex/build_artifacts.py}. Sau đó biên dịch bằng XeLaTeX/BibTeX hoặc Tectonic theo README. Đây chỉ là bước đọc dữ liệu và dàn trang, không huấn luyện mô hình, encode lại corpus hoặc gọi Gemini.
\begingroup\setstretch{1.0}\footnotesize\setlength{\tabcolsep}{4pt}\renewcommand{\arraystretch}{1.10}
\begin{longtable}{p{.29\linewidth}p{.62\linewidth}}
\caption{Bản đồ bằng chứng dùng để kiểm tra nội dung.}\label{tab:evidence}\\
\toprule Nội dung & Artifact \\
\midrule\endfirsthead
\multicolumn{2}{l}{\textit{Bảng tiếp theo}}\\
\toprule Nội dung & Artifact \\
\midrule\endhead
\bottomrule\endfoot
Corpus và disposition & \artifact{data/processed/main_corpus.parquet} \\
Số liệu xử lý & \artifact{outputs/reports/cleaning_summary.csv} \\
LDA grid / lựa chọn & \artifact{outputs/tables/stage2/lda_comparison.csv} \\
Proxy trước/sau và CI & \artifact{outputs/research_validation/proxy_test_metrics_ci.csv} \\
Coverage chunking & \artifact{outputs/bertopic_chunked/chunking_success.json} \\
Topic metrics & \artifact{outputs/bertopic_chunked/metrics.json} \\
Chuyển topic/noise và ARI & \artifact{outputs/research_validation/between_versions_agreement.json} \\
Rankings 14 methods & \artifact{outputs/search_upgrade_20261018/evaluation_14_methods/top5_fourteen_methods.csv} \\
Qrels AI và provenance & \artifact{outputs/retrieval_ai_evaluation_v1/ai_consensus_qrels.csv} \\
Agreement và coverage & \artifact{outputs/retrieval_ai_evaluation_v1/agreement.json} \\
Relevance AI một phần & \artifact{outputs/retrieval_ai_evaluation_v1/summary.csv} \\
Synthetic summary & \artifact{outputs/retrieval_ai_evaluation_v1/synthetic/summary.csv} \\
Tái tính 2.758 runs & \artifact{outputs/retrieval_ai_evaluation_v1/completion_validation.json} \\
Latency warm uncached & \artifact{outputs/search_upgrade_20261018/live_latency_scope.json} \\
237 tests đã thực hiện & \artifact{outputs/retrieval_ai_evaluation_v1/tests.xml} \\
Hash inputs AI task & \artifact{outputs/retrieval_ai_evaluation_v1/input_verification.json} \\
Parity baseline lịch sử & \artifact{outputs/retrieval_ai_evaluation_v1/legacy_parity.json} \\
Implementation lịch sử & \artifact{archive/cleanup_20261018/historical_files.zip} \\
\end{longtable}\endgroup
\source{Các đường dẫn tương đối từ workspace root}

\chapter{Figures thực nghiệm mở rộng}\label{app:figures}
\section{Thống kê mô tả corpus}
Các biểu đồ dưới đây là outputs đã có, không được tạo từ một corpus mới. Tần suất token, n-gram và metadata keywords phản ánh quy tắc xử lý tương ứng; không coi chúng là nhãn chủ đề chuyên gia. Các năm ít mẫu có thể gây biến động tỷ lệ lớn.
\evidencefigure{figures__stage2__eda_01_publication_trends.png}{Corpus đóng băng theo nguồn và năm. Các năm ít mẫu hoặc thiếu nguồn không đại diện toàn bộ hoạt động nghiên cứu.}{fig:output-18}{outputs/figures/stage2/eda_01_publication_trends.png}{.43}
\evidencefigure{figures__stage2__eda_02_top_terms.png}{Top token Stage 1 theo document frequency ở từng nguồn. Token phổ biến không đồng nghĩa chủ đề khoa học được chuyên gia xác nhận.}{fig:output-19}{outputs/figures/stage2/eda_02_top_terms.png}{.43}
\evidencefigure{figures__stage2__eda_03_ngrams.png}{Các n-gram liền kề trong chuỗi token Stage 1. Đây là thống kê sau preprocessing, không phải quan hệ hóa học đã xác minh.}{fig:output-20}{outputs/figures/stage2/eda_03_ngrams.png}{.43}
\evidencefigure{figures__stage2__eda_04_author_keywords.png}{Author keywords từ metadata Publication. Giữ cách viết nguồn và không trộn chúng vào labels đánh giá retrieval.}{fig:output-21}{outputs/figures/stage2/eda_04_author_keywords.png}{.43}
\evidencefigure{figures__stage2__eda_05_length_distributions.png}{Phân bố độ dài token LDA và ký tự Abstract. Density được chuẩn hóa theo nguồn; giá trị 0 không hiển thị trên trục log.}{fig:output-22}{outputs/figures/stage2/eda_05_length_distributions.png}{.43}
\evidencefigure{figures__stage2__eda_06_domain_term_trends.png}{Tỷ lệ hiện diện các thuật ngữ miền theo nguồn/năm. Dao động ở năm ít tài liệu cần được diễn giải thận trọng.}{fig:output-23}{outputs/figures/stage2/eda_06_domain_term_trends.png}{.43}
\FloatBarrier

\section{LDA theo thời gian và từ khóa}
Soft topic mass được thống kê theo năm và nguồn. Không dùng sự tăng tỷ lệ trong corpus để khẳng định tăng trưởng toàn bộ lĩnh vực, vì phạm vi thu thập và mức độ đầy đủ các năm không đồng nhất.
\evidencefigure{figures__stage2__lda_04_topic_year_heatmap.png}{Topic LDA theo năm trong từng nguồn, dựa soft mass. Ô trắng tương ứng không có documents, không phải bằng chứng chủ đề không tồn tại.}{fig:output-24}{outputs/figures/stage2/lda_04_topic_year_heatmap.png}{.43}
\evidencefigure{figures__stage2__lda_05_largest_topic_trends.png}{Xu hướng các topic LDA lớn theo soft mass; khác nhau về phạm vi năm giữa hai nguồn nên không so trực tiếp độ cao đường.}{fig:output-25}{outputs/figures/stage2/lda_05_largest_topic_trends.png}{.43}
\evidencefigure{figures__stage2__lda_06_keyword_heatmap.png}{Trọng số Top10 keywords của LDA K=20. Chỉ một tập từ hiển thị theo tổng trọng số; hàng không phải xác suất relevance.}{fig:output-26}{outputs/figures/stage2/lda_06_keyword_heatmap.png}{.43}
\FloatBarrier

\section{BERTopic baseline và toàn bộ keywords chunked}
Các hình baseline dùng 92 topics, trong khi bảy trang keywords chunked dùng 84 topics. Topic IDs của hai phiên bản không có ánh xạ đồng nhất mặc định. Từ khóa là c-TF-IDF trên tokens đóng băng, còn giữ các giới hạn preprocessing đã nêu.
\evidencefigure{figures__stage3__bertopic_01_distribution.png}{Phân bố baseline BERTopic 92 topics và noise -1, dùng embeddings 128 tokens; không thay thế phân bố chunked 84 topics.}{fig:output-27}{outputs/figures/stage3/bertopic_01_distribution.png}{.43}
\evidencefigure{figures__stage3__bertopic_02_keywords.png}{Keywords của các topic baseline được chọn trong artifact. Các trang này không bao phủ toàn bộ 92 topics và chưa được chuyên gia đặt tên.}{fig:output-28}{outputs/figures/stage3/bertopic_02_keywords.png}{.43}
\evidencefigure{figures__stage3__bertopic_03_similarity_heatmap.png}{Cosine giữa centroid embeddings của topic baseline. Mức tương đồng cao không chứng minh hai topic tương đương về ý nghĩa khoa học.}{fig:output-29}{outputs/figures/stage3/bertopic_03_similarity_heatmap.png}{.43}
\FloatBarrier

\evidencefigure{bertopic_chunked__figures__top10_keywords_01.png}{Top10 keywords BERTopic chunked, trang 1/7; bảy trang bao phủ 84 topics không phải noise. Trọng số c-TF-IDF không phải relevance grade.}{fig:output-35}{outputs/bertopic_chunked/figures/top10_keywords_01.png}{.83}
\evidencefigure{bertopic_chunked__figures__top10_keywords_02.png}{Top10 keywords BERTopic chunked, trang 2/7; bảy trang bao phủ 84 topics không phải noise. Trọng số c-TF-IDF không phải relevance grade.}{fig:output-36}{outputs/bertopic_chunked/figures/top10_keywords_02.png}{.83}
\evidencefigure{bertopic_chunked__figures__top10_keywords_03.png}{Top10 keywords BERTopic chunked, trang 3/7; bảy trang bao phủ 84 topics không phải noise. Trọng số c-TF-IDF không phải relevance grade.}{fig:output-37}{outputs/bertopic_chunked/figures/top10_keywords_03.png}{.83}
\evidencefigure{bertopic_chunked__figures__top10_keywords_04.png}{Top10 keywords BERTopic chunked, trang 4/7; bảy trang bao phủ 84 topics không phải noise. Trọng số c-TF-IDF không phải relevance grade.}{fig:output-38}{outputs/bertopic_chunked/figures/top10_keywords_04.png}{.83}
\evidencefigure{bertopic_chunked__figures__top10_keywords_05.png}{Top10 keywords BERTopic chunked, trang 5/7; bảy trang bao phủ 84 topics không phải noise. Trọng số c-TF-IDF không phải relevance grade.}{fig:output-39}{outputs/bertopic_chunked/figures/top10_keywords_05.png}{.83}
\evidencefigure{bertopic_chunked__figures__top10_keywords_06.png}{Top10 keywords BERTopic chunked, trang 6/7; bảy trang bao phủ 84 topics không phải noise. Trọng số c-TF-IDF không phải relevance grade.}{fig:output-40}{outputs/bertopic_chunked/figures/top10_keywords_06.png}{.83}
\evidencefigure{bertopic_chunked__figures__top10_keywords_07.png}{Top10 keywords BERTopic chunked, trang 7/7; bảy trang bao phủ 84 topics không phải noise. Trọng số c-TF-IDF không phải relevance grade.}{fig:output-41}{outputs/bertopic_chunked/figures/top10_keywords_07.png}{.83}
\FloatBarrier

\section{Backend tìm kiếm lịch sử và giao diện Stage5}
Các figures Stage4/Stage5 dưới đây được giữ như bằng chứng phiên bản lịch sử. Chúng không mô tả latency, giao diện hoặc mặc định của Search v2 hiện tại. Overlap và số lượt thay đổi ranking không phải thước đo relevance; ảnh giao diện không phải một thí nghiệm chất lượng độc lập.
\evidencefigure{figures__stage4__latency_components.png}{Latency thành phần của backend lịch sử: retrieval/ranking và query encoding đo riêng. Không phải latency Search v2 có cross-encoder.}{fig:output-30}{outputs/figures/stage4/latency_components.png}{.43}
\evidencefigure{figures__stage4__latency_distribution.png}{Phân bố warm latency backend lịch sử trên 24 queries, ba lượt/query/method, model nạp sẵn. Không gộp với sáu requests live của Search v2.}{fig:output-31}{outputs/figures/stage4/latency_distribution.png}{.43}
\evidencefigure{figures__stage4__result_composition.png}{Thành phần nguồn và tài liệu noise trong kết quả lịch sử. Đây là thống kê kết quả, không phải relevance, bias đã xác nhận hoặc tỷ lệ lỗi.}{fig:output-32}{outputs/figures/stage4/result_composition.png}{.43}
\evidencefigure{figures__stage4__se2_reranking_changes.png}{Số queries lịch sử thay đổi thứ tự/tập Top5 khi cộng topic bonus vào dense. Có thay đổi ranking không chứng minh chất lượng được cải thiện.}{fig:output-33}{outputs/figures/stage4/se2_reranking_changes.png}{.43}
\evidencefigure{figures__stage4__top5_overlap_heatmap.png}{Overlap Jaccard Top5 giữa các phương pháp lịch sử. Đồng thuận tài liệu không phải agreement relevance hoặc kiểm định tương đương.}{fig:output-34}{outputs/figures/stage4/top5_overlap_heatmap.png}{.43}
\FloatBarrier

\evidencefigure{figures__stage5__lda_dashboard.png}{Ảnh giao diện dashboard LDA ở phiên bản Stage5 lịch sử. Các số liệu và điều khiển là snapshot của phiên bản đó.}{fig:output-42}{outputs/figures/stage5/lda_dashboard.png}{.83}
\evidencefigure{figures__stage5__bertopic_dashboard.png}{Ảnh dashboard BERTopic baseline lịch sử, hiển thị 92 topics và noise baseline; không phải dashboard chunked hiện tại.}{fig:output-43}{outputs/figures/stage5/bertopic_dashboard.png}{.83}
\evidencefigure{figures__stage5__initial_search_form.png}{Ảnh giao diện tìm kiếm Stage5 lịch sử dùng SE1/SE2/SE3. Search v2 hiện tại có backend và mặc định khác.}{fig:output-44}{outputs/figures/stage5/initial_search_form.png}{.83}
\FloatBarrier

\section{Biểu đồ relevance AI chưa đủ mẫu}
Hình này chỉ có một query remaining đủ pool, khác mean hai query trình bày trong chương thực nghiệm. Việc đưa hình vào phụ lục nhằm minh bạch outputs hiện có, không bổ sung bằng chứng cho một tuyên bố thắng trên benchmark đầy đủ.
\evidencefigure{retrieval_ai_evaluation_v1__ai_relevance_metrics_remaining.png}{Relevance AI trên đúng một query remaining đủ pool (Q11, 1/18), khác bảng mean hai queries Q11/Q19. Không chọn phương pháp thắng từ biểu đồ partial này.}{fig:output-45}{outputs/retrieval_ai_evaluation_v1/ai_relevance_metrics_remaining.png}{.43}
\FloatBarrier

